\documentclass[11pt]{article}
\usepackage[T1]{fontenc}
\usepackage[utf8]{inputenc}
\usepackage{lmodern}
\usepackage[margin=1in,headheight=15pt]{geometry}
\usepackage{microtype}
\usepackage{amsmath,amssymb,amsthm,mathtools}
\usepackage{booktabs,array,longtable}
\usepackage{xcolor}
\usepackage{enumitem}
\usepackage{fancyhdr}
\usepackage{hyperref}
\definecolor{linkblue}{RGB}{27,65,99}
\hypersetup{colorlinks=true,linkcolor=linkblue,citecolor=linkblue,urlcolor=linkblue,
 pdftitle={A hexagonal-lattice proof of the triangular-array conjecture A195806},
 pdfauthor={Research report prepared with ChatGPT},
 pdfsubject={Exact enumeration, rational generating functions, and a two-bound refinement}}
\urlstyle{same}
\pagestyle{fancy}
\fancyhf{}
\fancyhead[L]{\small A hexagonal-lattice proof of A195806}
\fancyhead[R]{\small September 2026}
\fancyfoot[C]{\thepage}
\renewcommand{\headrulewidth}{0.3pt}
\setlength{\parskip}{0.35em}
\setlength{\parindent}{1.2em}
\setlength{\emergencystretch}{2em}
\numberwithin{equation}{section}
\newtheorem{theorem}{Theorem}[section]
\newtheorem{proposition}[theorem]{Proposition}
\newtheorem{lemma}[theorem]{Lemma}
\newtheorem{corollary}[theorem]{Corollary}
\theoremstyle{definition}
\newtheorem{definition}[theorem]{Definition}
\newtheorem{remark}[theorem]{Remark}
\newcommand{\Q}{\mathbb Q}
\newcommand{\Z}{\mathbb Z}
\newcommand{\R}{\mathbb R}
\newcommand{\N}{\mathbb Z_{\geq 0}}
\newcommand{\pos}[1]{\left(#1\right)_{+}}
\newcommand{\coef}[2]{\left[#1\right]#2}
\DeclareMathOperator{\Diag}{Diag}
\DeclareMathOperator{\vol}{vol}
\newcommand{\Aseq}{\textup{A195806}}

\title{\textbf{A hexagonal-lattice proof of the\newline triangular-array conjecture A195806}\\[0.5em]
\large Exact enumeration, a minimal recurrence, and a two-bound refinement}
\author{Research report prepared with ChatGPT}
\date{September 20, 2026}

\begin{document}
\maketitle
\begin{abstract}
We count triangular arrays of side length five with entries in
$\{0,\ldots,n\}$, zero corners, and equal sums on lines of equal length in the
three lattice directions. The explicit quasipolynomial and constant-coefficient
recurrence for this count are labeled conjectural in OEIS A195806 and originate
in work of Kauers and Koutschan. We give an elementary proof. An integral
parametrization reduces the twelve entries to two difference parameters and
four independent translations. Their admissible ranges are measured by two
hexagonal norms. A common dihedral symmetry reduces the count to a weighted
sum over a single chamber, with weights $1$, $6$, and $12$. Summing this formula
gives a reduced rational generating function and proves all six conjectured
residue-class formulas. The minimal quasiperiod is six, the minimal eventual
constant-coefficient recurrence has order fourteen, and the leading coefficient
is $65/648$. We also derive a bivariate generating function for separate bounds
on two classes of entries and compute its piecewise-polynomial leading volume.
The report distinguishes the correctly indexed OEIS formulas from an indexing
shift and a numerical typo in the printed conjecture. Reproducible exact
polynomial certificates and independent finite enumerations accompany the
proof.
\end{abstract}

\noindent\textbf{Research status.}
This manuscript supplies a complete mathematical argument for the explicit
OEIS formulas, rather than a conjecture inferred from additional terms. The
source search described below did not establish first-in-literature priority.
The manuscript has not been externally refereed or verified in a proof
assistant. Its computer checks are supplementary to the proof.

\clearpage
\tableofcontents
\clearpage

\section{The problem and the precise claim being proved}
\label{sec:problem}

\subsection{Definition and indexing}
Label the twelve noncorner entries of a triangular array as follows:
\begin{equation}\label{eq:array}
\begin{array}{ccccccccc}
&&&&0&&&&\\
&&&a&&b&&&\\
&&c&&d&&e&&\\
&f&&g&&h&&i&\\
0&&j&&k&&l&&0
\end{array}.
\end{equation}
The three directions are the horizontal rows and the two families of sloping
lines. For each fixed line length, the three line sums must agree. Lines of
\emph{different} lengths are not required to have the same sum. In the notation
of \eqref{eq:array}, the nontrivial equations are exactly
\begin{align}
a+b&=f+j=i+l,\label{eq:length2}\\
c+d+e&=c+g+k=e+h+k,\label{eq:length3}\\
f+g+h+i&=b+d+g+j=a+d+h+l,\label{eq:length4}\\
a+c+f&=j+k+l=b+e+i.\label{eq:length5}
\end{align}
The three lines of length one are the prescribed zero corners, so their
condition is automatic.

\begin{definition}
For an integer $n\geq0$, let $A(n)$ be the number of arrays
\eqref{eq:array} satisfying \eqref{eq:length2}--\eqref{eq:length5}, with all
twelve letters in $\{0,\ldots,n\}$.
\end{definition}

Thus $A(0)=1$. The OEIS entry has offset one and begins with $A(1)=16$;
we adjoin the natural value at zero to simplify generating functions.
The first few counts are
\[
A(0),A(1),\ldots,A(7)=1,16,105,496,1759,5052,12469,27412.
\]
This distinction between the entry bound and the position in a displayed
sequence will matter when comparing formulas.

\subsection{Why this is a documented conjecture}
R.\ H.\ Hardin introduced the sequence in 2011. As retrieved on September 20,
2026, its OEIS entry explicitly labels both a recurrence and a six-part
quasipolynomial as conjectured, attributing them to Manuel Kauers and
Christoph Koutschan \cite{oeis}. Their 2023 paper includes the corresponding
problem in Section 5.1 and an explicit formula as Conjecture 11
\cite[pp.~24--25]{kk}.

The 2023 paper already establishes that the sequence is a quasipolynomial and
hence is D-finite. The unresolved point identified there is the explicit
formula, not the bare existence of some recurrence. Accordingly, the result
of this report is a proof of the particular OEIS quasipolynomial and its
constant-coefficient recurrence, together with structural refinements.
Targeted searches for the sequence identifier and proof-related terms did not
locate a later resolution; this is evidence about the sources inspected, not
an exhaustive priority claim.

There is a small but important source issue. The displayed expression in
Conjecture 11 uses a one-step shift relative to our entry-bound indexing, and
one of its constants has a numerical typo. Section~\ref{sec:source} gives an
exact audit. Our principal target is the correctly indexed collection of
formulas on the OEIS page, not the erroneous literal printed expression.

\subsection{The route to a proof}
A direct enumeration examines $(n+1)^{12}$ fillings. General lattice-point
methods can in principle compute the answer, but the special linear system
has a much smaller description. We first solve it over the integers, then
count four independent intervals of translations. The remaining two-dimensional
sum has a symmetry common to both interval lengths. Restricting to a symmetry
chamber turns all maxima into linear functions. The generating function then
requires only geometric series and two differentiations.

This order of operations is essential. Interpolating six degree-six
polynomials from numerical values would not, by itself, prove that they count
the arrays. Here the all-index counting formula is established before any
coefficient extraction.

\section{Main enumeration theorem}
\label{sec:main}

\begin{theorem}\label{thm:main}
The ordinary generating function $F(z)=\sum_{n\geq0}A(n)z^n$ is
\begin{equation}\label{eq:mainGF}
 F(z)=\frac{P(z)}{(1-z)^7(1+z)(1+z+z^2)^3}
     =\frac{P(z)}{(1-z)^3(1-z^2)(1-z^3)^3},
\end{equation}
where
\begin{align}\label{eq:P}
P(z)={}&1+13z+59z^2+212z^3+471z^4+753z^5+882z^6\notag\\
      &\quad+753z^7+471z^8+212z^9+59z^{10}+13z^{11}+z^{12}.
\end{align}
For $n=6k+r$, where $k\geq0$ and $0\leq r\leq5$, write
$A(6k+r)=\sum_{d=0}^6 c_{r,d}k^d$. The coefficients are those in
Table~\ref{tab:residues}. In particular, the conjectured OEIS quasipolynomial
is correct for every index in its stated domain.
\end{theorem}

\begin{table}[htbp]
\centering
\small
\setlength{\tabcolsep}{5pt}
\begin{tabular}{crrrrrrr}
\toprule
$r$&$c_{r,6}$&$c_{r,5}$&$c_{r,4}$&$c_{r,3}$&$c_{r,2}$&$c_{r,1}$&$c_{r,0}$\\
\midrule
0&4680&4680&2275&650&158&25&1\\
1&4680&9360&8125&3900&1133&198&16\\
2&4680&14040&17875&12350&4922&1087&105\\
3&4680&18720&31525&28600&14783&4148&496\\
4&4680&23400&49075&55250&35258&12121&1759\\
5&4680&28080&70525&94900&72197&29474&5052\\
\bottomrule
\end{tabular}
\caption{Exact coefficients of $A(6k+r)$, with powers in descending order.
These agree with the conjectured OEIS formulas \cite{oeis}; the final row is
written in expanded form.}
\label{tab:residues}
\end{table}

For example, the final row can equivalently be written in the factored form
\begin{equation}\label{eq:lastresidue}
 A(6k+5)=(k+1)^2
 \bigl(4680k^4+18720k^3+28405k^2+19370k+5052\bigr).
\end{equation}
The proof of Theorem~\ref{thm:main} occupies Sections~\ref{sec:param}
through~\ref{sec:extract}. The period and recurrence are treated afterward,
so that their minimality does not have to be assumed in deriving the formula.

\section{An integral six-parameter description}
\label{sec:param}

\begin{lemma}\label{lem:param}
All integer solutions of \eqref{eq:length2}--\eqref{eq:length5}, without entry
bounds, are given uniquely by six integers
$(x,y,\alpha,\beta,\gamma,\delta)$ through
\begin{equation}\label{eq:paramarray}
\setlength{\arraycolsep}{3pt}
\begin{array}{ccccccccc}
&&&&0&&&&\\
&&&\beta+y&&\alpha+x-y&&&\\
&&\delta+2x-y&&\gamma&&\delta+x+y&&\\
&\alpha&&\gamma+x+y&&\gamma+2x-y&&\beta&\\
0&&\beta+x&&\delta&&\alpha+x&&0
\end{array}.
\end{equation}
The inverse map is
\begin{equation}\label{eq:inverse}
x=j-i,\qquad y=a-i,\qquad
\alpha=f,\quad\beta=i,\quad\gamma=d,\quad\delta=k.
\end{equation}
\end{lemma}

\begin{proof}
Set $\alpha=f$, $\beta=i$, $\gamma=d$, and $\delta=k$. From
$f+j=i+l$ we obtain
\[
 j-i=l-f=:x,
\]
so $j=\beta+x$ and $l=\alpha+x$. Set $y=a-i$; then
$a=\beta+y$. The remaining length-two equality determines
\[
 b=f+j-a=\alpha+x-y.
\]
Next, the common length-five sum is
\[
 j+k+l=\alpha+\beta+\delta+2x.
\]
Equating this with $a+c+f$ and $b+e+i$ gives, respectively,
\[
 c=\delta+2x-y,\qquad e=\delta+x+y.
\]
The length-three equations now yield
\[
 g=d+e-k=\gamma+x+y,\qquad
 h=c+d-k=\gamma+2x-y.
\]
Every entry has therefore been forced into the form
\eqref{eq:paramarray}. This proves necessity and uniqueness.

Conversely, substitute arbitrary integers into \eqref{eq:paramarray}. The
three sums within each length class become, respectively,
\[
\begin{array}{c|c}
\text{line length}&\text{common sum}\\ \hline
2&\alpha+\beta+x\\
3&2\delta+\gamma+3x\\
4&\alpha+\beta+2\gamma+3x\\
5&\alpha+\beta+\delta+2x.
\end{array}
\]
All required equalities hold, including the length-four equations not used in
the elimination. The inverse formulas are integral, so no congruence
conditions have been lost.
\end{proof}

\begin{remark}
The unconstrained real solution space has dimension six, and its integer
lattice is identified with $\Z^6$ by the same parametrization. The dimension
statement alone would not suffice for counting: a merely rational
parametrization could hide divisibility restrictions. The explicit integral
inverse in \eqref{eq:inverse} is the point that rules them out.
\end{remark}

\section{Translation intervals and two hexagonal norms}
\label{sec:norms}

For a real number $t$, write $t_+=\max(t,0)$.

\begin{lemma}\label{lem:range}
For integers $u,v$ and $n\geq0$, the number of integers $t$ for which
$t,t+u,t+v$ all belong to $\{0,\ldots,n\}$ is
\begin{equation}\label{eq:triplecount}
\bigl(n+1-\max(|u|,|v|,|u-v|)\bigr)_+.
\end{equation}
\end{lemma}

\begin{proof}
The allowed interval for $t$ is
\[
 -\min(0,u,v)\ \leq t\ \leq n-\max(0,u,v).
\]
Its number of integer points is the positive part of its upper endpoint minus
its lower endpoint plus one. The difference between the largest and smallest
of three real numbers is the largest of their three pairwise distances.
Applied to $0,u,v$, this gives \eqref{eq:triplecount}.
\end{proof}

Define
\begin{align}
 H_1(x,y)&=\max\{|x|,|y|,|x-y|\},\label{eq:H1}\\
 H_2(x,y)&=\max\{|x+y|,|2x-y|,|x-2y|\}.\label{eq:H2}
\end{align}
Both are norms on $\R^2$ with centrally symmetric hexagonal unit balls.
Only their explicit maximum formulas are needed below.

For fixed $x,y$, the four translation parameters in
\eqref{eq:paramarray} occur in four disjoint triples. Their offsets are
\begin{equation}\label{eq:offsets}
\begin{array}{c|c|c}
\text{translation}&\text{offsets}&\text{range}\\ \hline
\beta&0,y,x&H_1(x,y)\\
\alpha&0,x-y,x&H_1(x,y)\\
\gamma&0,x+y,2x-y&H_2(x,y)\\
\delta&0,2x-y,x+y&H_2(x,y).
\end{array}
\end{equation}
The first two ranges agree because their pairwise differences, up to sign,
are $x,y,x-y$. The last two triples simply have their nonzero offsets
interchanged.

\begin{proposition}\label{prop:normcount}
For every $n\geq0$,
\begin{equation}\label{eq:normcount}
 A(n)=\sum_{(x,y)\in\Z^2}
 \bigl(n+1-H_1(x,y)\bigr)_+^2
 \bigl(n+1-H_2(x,y)\bigr)_+^2.
\end{equation}
\end{proposition}

\begin{proof}
The parametrization is bijective, and the four translation parameters are
independent once $x,y$ have been fixed. Lemma~\ref{lem:range} therefore gives
the product in each summand. A summand is nonzero only if $H_1(x,y)\leq n$,
which in particular implies $|x|,|y|\leq n$. Thus the displayed sum has only
finitely many nonzero terms and accounts for every bounded array exactly
once.
\end{proof}

This already gives an $O(n^2)$-term counting formula, replacing the original
twelve-dimensional search. The remaining maxima are eliminated by exploiting
a symmetry shared by both norms.

\section{A common dihedral chamber}
\label{sec:chamber}

\begin{lemma}\label{lem:dihedral}
The transformations
\begin{equation}\label{eq:TS}
 T(x,y)=(x-y,x),\qquad S(x,y)=(x,x-y)
\end{equation}
generate a group of order twelve preserving $\Z^2$, $H_1$, and $H_2$.
Each orbit meets the closed chamber
\begin{equation}\label{eq:chamber}
\mathcal C=\{(x,y)\in\R^2:x\geq2y\geq0\}
\end{equation}
in exactly one point. A nonzero chamber point has orbit size six on either
boundary ray and twelve in the interior; the origin has orbit size one.
\end{lemma}

\begin{proof}
A direct computation gives $T^6=S^2=1$ and $STS=T^{-1}$.
For example, the three linear forms defining $H_2$ are transformed by $T$
into
\[
 2x-y,\quad x-2y,\quad -(x+y),
\]
so their absolute values are permuted. The same check works for $S$ and for
the defining forms of $H_1$.

For completeness, the fundamental-domain assertion can be seen in the
positive-definite quadratic form
\[
 q_0(x,y)=x^2-xy+y^2.
\]
Both transformations preserve this form. In the corresponding Euclidean
metric, $T$ is a rotation through $60$ degrees: its determinant is one and
its trace is one. The map $S$ is the reflection fixing the line $x=2y$.
The reflection $T^{-1}S$ fixes the line $y=0$. The positive rays on these two
lines enclose an angle of $30$ degrees. Their images divide the plane into
twelve sectors, whose interiors are disjoint. This proves that the group has
order twelve and that the closed sector \eqref{eq:chamber} contains exactly
one representative of each orbit.

An interior nonzero point is fixed by no nonidentity group element. A point
on either boundary ray is fixed by exactly one reflection, so its stabilizer
has order two. The orbit sizes follow, while every group element fixes the
origin.
\end{proof}

\begin{remark}
This is a symmetry of the auxiliary parameter lattice. We are \emph{not}
counting arrays modulo rotations or reflections. All arrays remain labeled
by their positions, and the orbit sizes are multiplicities in a sum, not
identifications of arrays.
\end{remark}

Use integral coordinates on the chamber by setting
\begin{equation}\label{eq:pq}
 p=x-2y,\qquad q=y,
 \qquad\text{so that}\qquad x=p+2q,\quad y=q.
\end{equation}
This is a bijection between $\mathcal C\cap\Z^2$ and $\N^2$. On the chamber,
\begin{equation}\label{eq:linearH}
 H_1=p+2q,\qquad H_2=2p+3q.
\end{equation}
Indeed, $x$ is the largest defining absolute value for $H_1$, and $2x-y$
is the largest for $H_2$. In particular,
\begin{equation}\label{eq:compareH}
 \tfrac32H_1\leq H_2\leq2H_1
\end{equation}
holds globally, by symmetry and \eqref{eq:linearH}.

Define the multiplicity
\begin{equation}\label{eq:mu}
 \mu(p,q)=
 \begin{cases}
 1,&p=q=0,\\
 6,&\text{exactly one of }p,q\text{ is zero},\\
 12,&p,q>0.
 \end{cases}
\end{equation}
The equal-bound count can now be written with no maxima.

\begin{theorem}\label{thm:sum}
For every integer $n\geq0$,
\begin{equation}\label{eq:positivesum}
 \boxed{\displaystyle
 A(n)=\sum_{\substack{p,q,r\geq0\\2p+3q+r=n}}
 \mu(p,q)(r+1)^2(r+p+q+1)^2.}
\end{equation}
\end{theorem}

\begin{proof}
Apply Lemma~\ref{lem:dihedral} to \eqref{eq:normcount}. Since
$H_2\geq H_1$, a nonzero chamber summand is characterized by
$2p+3q\leq n$. Set $r=n-2p-3q$. Then
\[
 n+1-H_2=r+1,\qquad n+1-H_1=r+p+q+1.
\]
The orbit multiplicities give exactly \eqref{eq:positivesum}.
\end{proof}

The first values illustrate the boundary weights. For $n=1$, only $p=q=0$
occurs, so $A(1)=2^4=16$. For $n=2$, one new point occurs on the $q=0$
boundary:
\[
 A(2)=3^4+6\cdot1^2\cdot2^2=105.
\]
For $n=3$, the two nonzero possibilities are $(p,q,r)=(1,0,1)$ and
$(0,1,0)$, giving
\[
 A(3)=4^4+6\cdot2^2\cdot3^2+6\cdot1^2\cdot2^2=496.
\]
These examples are checks on the formula; its proof has no finite-index
restriction.

\section{Summing the chamber formula}
\label{sec:GF}

Let
\begin{equation}\label{eq:G}
 G(u,v)=\sum_{p,q\geq0}\mu(p,q)u^pv^q.
\end{equation}
Separate the origin, the two axes, and the open quadrant to obtain
\begin{align}\label{eq:Grational}
 G(u,v)
 &=1+\frac{6u}{1-u}+\frac{6v}{1-v}
       +\frac{12uv}{(1-u)(1-v)}\notag\\
 &=\frac{1+5u+5v+uv}{(1-u)(1-v)}.
\end{align}
The Euler operator
\[
 D=u\frac{\partial}{\partial u}+v\frac{\partial}{\partial v}
\]
multiplies the coefficient of $u^pv^q$ by $p+q$. Hence $DG$ and $D^2G$
encode the first two powers of that sum. Direct differentiation gives
\begin{align}
 DG&=\frac{6(u+v)(1-uv)}{(1-u)^2(1-v)^2},\label{eq:DG}\\
 D^2G&=\frac{6(u+v)(1+u+v-6uv+u^2v+uv^2+u^2v^2)}
 {(1-u)^3(1-v)^3}.\label{eq:D2G}
\end{align}

For $j\geq0$, put $S_j(z)=\sum_{r\geq0}(r+1)^jz^r$. The needed cases are
\begin{equation}\label{eq:S}
\begin{split}
 S_2(z)&=\frac{1+z}{(1-z)^3},\qquad
 S_3(z)=\frac{1+4z+z^2}{(1-z)^4},\\
 S_4(z)&=\frac{1+11z+11z^2+z^3}{(1-z)^5}.
\end{split}
\end{equation}
These follow from the geometric series by applying $z\,d/dz+1$ repeatedly.
The weight in \eqref{eq:positivesum} expands as
\[
 (r+1)^2(r+p+q+1)^2
 =(r+1)^4+2(p+q)(r+1)^3+(p+q)^2(r+1)^2.
\]
Multiplying by $z^{2p+3q+r}$ and summing therefore yields the exact formal
power-series identity
\begin{equation}\label{eq:GFoperator}
 F(z)=\left.
 \bigl(S_4(z)G(u,v)+2S_3(z)DG(u,v)+S_2(z)D^2G(u,v)\bigr)
 \right|_{u=z^2,\ v=z^3}.
\end{equation}
Each coefficient involves finitely many triples, so there is no analytic
interchange-of-limits issue.

\subsection{An explicit certificate for the rational simplification}
To make the calculation from \eqref{eq:GFoperator} to \eqref{eq:mainGF}
checkable without trusting a computer algebra simplifier, set $U=z^2$,
$V=z^3$, and define
\begin{align*}
 M_0&=1+5U+5V+UV,\\
 M_1&=6(U+V)(1-UV),\\
 M_2&=6(U+V)(1+U+V-6UV+U^2V+UV^2+U^2V^2).
\end{align*}
The common numerator on the right side of \eqref{eq:GFoperator}, over the
denominator $(1-z)^5(1-U)^3(1-V)^3$, is
\begin{align}\label{eq:polycertificate}
 &(1+11z+11z^2+z^3)M_0(1-U)^2(1-V)^2\notag\\
 &\quad+2(1+4z+z^2)(1-z)M_1(1-U)(1-V)\notag\\
 &\quad+(1+z)(1-z)^2M_2
   =(1-z)^2(1-U)^2P(z).
\end{align}
Equation~\eqref{eq:polycertificate} is an identity between integer
polynomials of degree at most eighteen. It is verified simply by expanding
and comparing coefficients; the supplied standard-library verifier performs
exactly that multiplication and comparison.

Canceling the displayed factors gives
\[
 F(z)=\frac{P(z)}{(1-z)^3(1-U)(1-V)^3}.
\]
Substitute $U=z^2$, $V=z^3$ and factor $1-z^2$ and $1-z^3$ to obtain
\eqref{eq:mainGF}. This establishes the generating-function part of
Theorem~\ref{thm:main} directly from the combinatorial definition.

\section{Exact coefficient extraction}
\label{sec:extract}

\subsection{Polynomial and periodic parts}
A compact form of the answer separates a polynomial in the centered variable
$N=n+1$ from periodic corrections. Define
\begin{equation}\label{eq:M}
 M(N)=\frac{65N^6}{648}+\frac{325N^4}{1296}
      +\frac{439N^2}{324}-\frac{2521}{7776}
\end{equation}
and
\begin{equation}\label{eq:B}
\begin{split}
 B_0(N)&=\frac{2(3N^2-24N-14)}{243},\\
 B_1(N)&=\frac{2(3N^2+24N-14)}{243},\\
 B_2(N)&=-\frac{4(3N^2-14)}{243}.
\end{split}
\end{equation}

\begin{proposition}\label{prop:centered}
For every $n\geq0$,
\begin{equation}\label{eq:centered}
 \boxed{\displaystyle
 A(n)=M(n+1)-\frac{3}{32}(-1)^n+B_{n\bmod3}(n+1).}
\end{equation}
\end{proposition}

\begin{proof}
Write $C=1+z+z^2$. Partial fractions of \eqref{eq:mainGF} give
\begin{align}\label{eq:PF}
 F(z)={}&\frac{650}{9(1-z)^7}-\frac{1625}{9(1-z)^6}
       +\frac{325}{2(1-z)^5}-\frac{2275}{36(1-z)^4}\notag\\
 &+\frac{2687}{216(1-z)^3}-\frac{737}{432(1-z)^2}
       -\frac{2521}{7776(1-z)}-\frac{3}{32(1+z)}\notag\\
 &-\frac{14(4z-1)}{243C}
       -\frac{2(13z+8)}{81C^2}
       +\frac{4(z-1)}{27C^3}.
\end{align}
This identity can be checked by multiplying by
$7776(1-z)^7(1+z)C^3$; the resulting integer-polynomial equality is a
second directly checkable algebra calculation, included in the verifier.

The identity
\[
 [z^n](1-z)^{-j}=\binom{n+j-1}{j-1}
\]
shows that the first seven terms in \eqref{eq:PF} have coefficient $M(n+1)$.
The $(1+z)^{-1}$ term contributes $-3(-1)^n/32$.

For the last three terms, use
\begin{equation}\label{eq:Cexpand}
 C^{-j}=\frac{(1-z)^j}{(1-z^3)^j},\qquad
 (1-z^3)^{-j}=\sum_{k\geq0}\binom{k+j-1}{j-1}z^{3k}.
\end{equation}
Expanding the short numerator $(1-z)^j$ and separating the three residue
classes gives the following coefficients for their sum:
\[
\begin{array}{c|c}
n&\text{coefficient of the last three terms of \eqref{eq:PF}}\\ \hline
3k&(54k^2-108k-70)/243\\
3k+1&(54k^2+216k+92)/243\\
3k+2&(-108k^2-216k-52)/243.
\end{array}
\]
These are exactly $B_{n\bmod3}(n+1)$ in \eqref{eq:B}. All extractions hold
also at $k=0$, either directly from \eqref{eq:Cexpand} or by substituting
into the three displayed polynomials. This proves \eqref{eq:centered}.
\end{proof}

Substitute $n=6k+r$ into \eqref{eq:centered}. The parity and residue modulo
three are then fixed, so ordinary polynomial expansion gives each row of
Table~\ref{tab:residues}. This completes the proof of
Theorem~\ref{thm:main}. Notice that the periodic corrections have degree at
most two. Consequently, the coefficients of $n^6,n^5,n^4,n^3$ do not depend
on the residue class.

\subsection{An independent all-index certificate for the residue table}
The table can also be verified directly against the rational function,
without using the partial fraction decomposition. Let polynomials $T_d(t)$
be defined by
\begin{equation}\label{eq:EulerT}
 T_0(t)=1,\qquad
 T_{d+1}(t)=t\bigl((1-t)T'_d(t)+(d+1)T_d(t)\bigr).
\end{equation}
Differentiating a geometric series shows that
\[
 \sum_{k\geq0}k^dt^k=\frac{T_d(t)}{(1-t)^{d+1}}
\]
for $d\geq1$; for $d=0$ the left side is understood as $\sum_{k\geq0}t^k$.
Define the integer polynomial
\begin{equation}\label{eq:Hcertificate}
 H(z)=\sum_{r=0}^5z^r\sum_{d=0}^6
 c_{r,d}T_d(z^6)(1-z^6)^{6-d}.
\end{equation}
The generating function of the residue table is $H(z)/(1-z^6)^7$.
Let
\begin{equation}\label{eq:Q}
 Q(z)=(1-z)^7(1+z)(1+z+z^2)^3.
\end{equation}
A third finite integer-polynomial identity is
\begin{equation}\label{eq:residuecertificate}
 H(z)Q(z)=P(z)(1-z^6)^7.
\end{equation}
The degree on both sides is at most fifty-four. The included program constructs
$T_d$, $H$, $P$, and $Q$ as coefficient lists and checks this identity exactly.
It certifies equality of the \emph{whole} generating functions, not merely
agreement of a finite prefix. Since their denominators have constant term one,
\eqref{eq:residuecertificate} proves coefficientwise equality in
$\Z[[z]]$.

\section{Minimal period, minimal recurrence, and growth}
\label{sec:consequences}

\begin{proposition}\label{prop:reduced}
The fraction in \eqref{eq:mainGF} is reduced over $\Q[z]$.
\end{proposition}

\begin{proof}
Direct evaluation and polynomial division give
\begin{equation}\label{eq:noncancel}
 P(1)=3900,\qquad P(-1)=-12,\qquad
 P(z)\equiv12\pmod{1+z+z^2}.
\end{equation}
These are all irreducible factors of the denominator, and none divides $P$.
\end{proof}

\begin{corollary}\label{cor:period}
The sequence $A(n)$ has degree six and minimal quasiperiod six.
\end{corollary}

\begin{proof}
Table~\ref{tab:residues} supplies quasiperiod six and degree six. Suppose a
quasipolynomial representation had quasiperiod $h$. Summing each of its
polynomial residue classes would express its generating function with poles
only at $h$th roots of unity: a common denominator is a power of $1-z^h$.
But the reduced generating function has a pole at $-1$ and at both primitive
cube roots of unity by \eqref{eq:noncancel}. Thus $2\mid h$ and $3\mid h$,
so $6\mid h$. No smaller positive quasiperiod is possible.
\end{proof}

Let $E$ denote the forward shift on sequences, $(EA)(n)=A(n+1)$.

\begin{corollary}\label{cor:recurrence}
For every $n\geq0$,
\begin{equation}\label{eq:operatorrec}
 (E-1)^7(E+1)(E^2+E+1)^3A(n)=0.
\end{equation}
The minimal order of an eventual homogeneous linear recurrence with constant
coefficients over $\Q$ is fourteen.
\end{corollary}

\begin{proof}
The generating function has denominator degree fourteen and numerator degree
twelve. Multiplication by $Q(z)$ gives $Q(z)F(z)=P(z)$, so every coefficient
of degree at least fourteen vanishes. Reversing the denominator gives the
forward-shift recurrence \eqref{eq:operatorrec}. In fact
$z^{14}Q(1/z)=-Q(z)$, so the sign change is immaterial.

Conversely, any eventual homogeneous constant-coefficient recurrence of order
$d$ yields a polynomial multiplier of degree at most $d$ that makes $F(z)$ a
polynomial: the recurrence cancels the entire sufficiently long tail. Since
$P/Q$ is reduced, that multiplier must be divisible by $Q$. Therefore
$d\geq\deg Q=14$.
\end{proof}

The qualifier \emph{constant coefficients} is important. This result is not
a minimality claim among recurrences with coefficients depending polynomially
on $n$. The 2023 source reports a smaller-order polynomial-coefficient guess;
it is a different notion of order \cite{kk}. The constant-coefficient
recurrence conjectured on the OEIS page is precisely the expanded form
\begin{align}\label{eq:expandedrec}
0={}&A(n)-3A(n+1)+2A(n+2)-A(n+3)+6A(n+4)\notag\\
 &-5A(n+5)-3A(n+6)+3A(n+8)+5A(n+9)\notag\\
 &-6A(n+10)+A(n+11)-2A(n+12)+3A(n+13)-A(n+14).
\end{align}
There is no $A(n+7)$ term.

\begin{corollary}\label{cor:growth}
As $n\to\infty$,
\begin{equation}\label{eq:asymptotic}
 A(n)=\frac{65}{648}n^6+\frac{65}{108}n^5
     +\frac{2275}{1296}n^4+\frac{325}{108}n^3+O(n^2).
\end{equation}
In particular, $A(n)\sim65n^6/648$.
\end{corollary}

\begin{proof}
Expand $M(n+1)$ in \eqref{eq:centered}; the remaining terms have size
$O(n^2)$. Independently, the leading constant follows from the pole at one:
\[
 \lim_{z\to1}(1-z)^7F(z)=\frac{P(1)}{2\cdot3^3}
 =\frac{650}{9},
\]
and the coefficient of $(1-z)^{-7}$ has leading term $n^6/6!$.
\end{proof}

The centered exact formula is more informative than an asymptotic expansion:
it records every lower-order oscillation, including the constant alternating
term. No exponentially small or uncomputed remainder is needed for exact
evaluation.

\section{Reciprocity and interior arrays}
\label{sec:reciprocity}

The numerator has the palindromic symmetry
$z^{12}P(1/z)=P(z)$, while the denominator is antipalindromic. Consequently,
\begin{equation}\label{eq:rationalreciprocity}
 F(1/z)=-z^2F(z)
\end{equation}
as an identity of rational functions. This is not a termwise identification
of the power-series expansions around zero and infinity.

The residue formulas uniquely extend $A(n)$ from nonnegative to all integer
arguments: use the polynomial assigned to the residue class of $n$. Denote
this continuation by $\widetilde A(n)$.

\begin{proposition}\label{prop:reciprocity}
For every integer $n$,
\begin{equation}\label{eq:quasireciprocity}
 \widetilde A(-n)=\widetilde A(n-2).
\end{equation}
For $n\geq2$, the number of valid arrays with all twelve noncorner entries
strictly between $0$ and $n$ is $A(n-2)$.
\end{proposition}

\begin{proof}
The polynomial $M(N)$ is even, and
\[
 B_0(-N)=B_1(N),\qquad B_1(-N)=B_0(N),\qquad B_2(-N)=B_2(N).
\]
Under the replacement of the sequence argument $t$ by $-t-2$, the centered
variable $N=t+1$ changes sign, the parity of $t$ stays fixed, and the residues
$0,1,2$ modulo three are changed to $1,0,2$, respectively. Formula
\eqref{eq:centered}, regarded as a quasipolynomial identity, therefore gives
$\widetilde A(t)=\widetilde A(-t-2)$. Substitute $t=n-2$.

For the counting assertion, subtract one from each noncorner entry and keep
the three corners zero. Within a fixed line-length class, every line contains
the same number of noncorner positions: those numbers are $2,3,4,3$ for
lengths $2,3,4,5$. Hence each equality of line sums is preserved. The strict
integer bounds $1\leq a,\ldots,l\leq n-1$ become the closed bounds
$0\leq a,\ldots,l\leq n-2$. Adding one gives the inverse bijection.
\end{proof}

The geometric interpretation is also exact. In the six coordinates of
Lemma~\ref{lem:param}, let $\mathcal P$ be the real polytope obtained by
bounding every noncorner entry between zero and one. The point
$x=y=0$, $\alpha=\beta=\gamma=\delta=1/2$ satisfies every bound strictly,
so $\mathcal P$ is six-dimensional. Its relative interior consists exactly
of points where every noncorner entry is strictly between the bounds.
Indeed, each coordinate-bound functional is nonconstant on the solution
space, and a point attaining such a bound lies on a supporting hyperplane;
conversely, finitely many strict linear inequalities remain strict in a
small neighborhood.

Thus the preceding bijection describes the interior lattice-point count
of $n\mathcal P$. At $n=1$ there are no such points and
$\widetilde A(-1)=0$. At $n=2$ there is exactly one: all twelve noncorner
entries are one. These statements and \eqref{eq:quasireciprocity} have been
proved directly, without needing a general reciprocity theorem as a black box.

\section{Two entry bounds and a bivariate generating function}
\label{sec:twobound}

The parametrization naturally divides the entries into two classes:
\begin{equation}\label{eq:entryclasses}
 \mathcal U=\{a,b,f,i,j,l\},\qquad
 \mathcal V=\{c,d,e,g,h,k\}.
\end{equation}
The first class contains the six positions adjacent to corners. The second
contains the three side midpoints and the three interior positions.
Let $A(n,m)$ count the same arrays, with entries in $\mathcal U$ bounded
between zero and $n$ and entries in $\mathcal V$ bounded between zero and $m$.
The corners remain zero, and $A(n,n)=A(n)$.

\begin{theorem}\label{thm:twobound}
For integers $n,m\geq0$,
\begin{equation}\label{eq:mixedsum}
 A(n,m)=\sum_{p,q\geq0}\mu(p,q)
 \bigl(n+1-p-2q\bigr)_+^2
 \bigl(m+1-2p-3q\bigr)_+^2.
\end{equation}
Moreover,
\begin{align}\label{eq:bivGF}
 &\sum_{n,m\geq0}A(n,m)X^nY^m\notag\\
 &\qquad=\frac{(1+X)(1+Y)}{(1-X)^3(1-Y)^3}
 \frac{1+5XY^2+5X^2Y^3+X^3Y^5}
 {(1-XY^2)(1-X^2Y^3)}.
\end{align}
\end{theorem}

\begin{proof}
In \eqref{eq:offsets}, the $\alpha$- and $\beta$-triples are exactly
$\mathcal U$, while the $\gamma$- and $\delta$-triples are exactly
$\mathcal V$. Their independent translation counts are therefore
$(n+1-H_1)_+^2(m+1-H_2)_+^2$. The same chamber argument proves
\eqref{eq:mixedsum}; no comparison between $n$ and $m$ is needed.

For a fixed chamber point, write
\[
 n=p+2q+r,\qquad m=2p+3q+s,\qquad r,s\geq0.
\]
Its contribution is $\mu(p,q)(r+1)^2(s+1)^2$. Consequently the bivariate
series factors as
\[
 S_2(X)S_2(Y)G(XY^2,X^2Y^3).
\]
Substitution of \eqref{eq:Grational} and \eqref{eq:S} gives
\eqref{eq:bivGF}.
\end{proof}

As boundary checks,
\[
 A(n,0)=(n+1)^2,\qquad A(0,m)=(m+1)^2.
\]
In either case the narrower bound forces $x=y=0$ and leaves two independent
translations for the other class.

The equal-bound generating function is the \emph{diagonal} of the bivariate
series: extract the coefficients of $X^nY^n$ and attach $z^n$.
It is not obtained by substituting $X=Y=z$, which would group terms by
$n+m$. This distinction explains why the very short bivariate rational form
can coexist with the more elaborate univariate numerator in \eqref{eq:P}.

\section{The leading volume under unequal scaling}
\label{sec:volume}

The two-bound formula also yields a continuous leading coefficient with an
explicit change of shape. For fixed real $a,b>0$, define
\begin{equation}\label{eq:volumeintegral}
 V(a,b)=12\iint_{\substack{p,q\geq0\\p+2q\leq a\\2p+3q\leq b}}
 (a-p-2q)^2(b-2p-3q)^2\,dp\,dq.
\end{equation}

\begin{theorem}\label{thm:volume}
For fixed $a,b>0$, as $t\to\infty$,
\begin{equation}\label{eq:anisotropicasymptotic}
 A(\lfloor at\rfloor,\lfloor bt\rfloor)
 =V(a,b)t^6+O_{a,b}(t^5).
\end{equation}
The leading coefficient has the truncated-power expression
\begin{equation}\label{eq:volumetruncated}
 \boxed{\displaystyle
 V(a,b)=\frac{b^4(540a^2-252ab+37b^2)
 -(2b-3a)_+^6+27(b-2a)_+^6}{3240}.}
\end{equation}
\end{theorem}

\begin{proof}[Proof of the asymptotic formula]
Rescale the chamber variables in \eqref{eq:mixedsum} by $t$. There are
$O_{a,b}(t^2)$ contributing lattice points. At an interior chamber point,
the multiplicity is twelve. Replacing each integer bound plus one by $at$
or $bt$ changes the summand by $O_{a,b}(t^3)$ on a common bounded rescaled
support, hence changes the total by $O_{a,b}(t^5)$.

The two axes contain only $O_{a,b}(t)$ relevant points. Replacing their
multiplicity six by twelve, and treating the origin similarly, changes the
sum by $O_{a,b}(t^5)$ because each weight is $O_{a,b}(t^4)$. The remaining
expression is $t^6$ times a two-dimensional Riemann sum for the compactly
supported function
\[
 12(a-p-2q)_+^2(b-2p-3q)_+^2
 \quad\text{on the first quadrant}.
\]
On a fixed rectangle containing its support this function is Lipschitz.
The Riemann-sum error is $O_{a,b}(1/t)$, obtained by bounding the variation
on each square of side $1/t$. Multiplying by $t^6$ gives the stated
$O_{a,b}(t^5)$ error and the integral \eqref{eq:volumeintegral}.
\end{proof}

\begin{proof}[Evaluation of the integral]
The two upper bounds on $p$ are
\[
 p\leq a-2q,\qquad p\leq\frac{b-3q}{2}.
\]
Their intersection is
\[
 (p,q)=(2b-3a,\,2a-b).
\]
The integration region changes combinatorial type at $b=3a/2$ and $b=2a$.

If $b\leq3a/2$, only the second sloping bound is active. Integration over
$0\leq q\leq b/3$ and $0\leq p\leq(b-3q)/2$ yields
\begin{equation}\label{eq:volsmall}
 V(a,b)=\frac{b^4(540a^2-252ab+37b^2)}{3240}.
\end{equation}
If $b\geq2a$, only the first sloping bound is active. Integration over
$0\leq q\leq a/2$ and $0\leq p\leq a-2q$ gives
\begin{equation}\label{eq:vollarge}
 V(a,b)=\frac{a^4(37a^2-84ab+60b^2)}{120}.
\end{equation}
In the intermediate region $3a/2\leq b\leq2a$, split at $q=2a-b$.
For $0\leq q\leq2a-b$, use the upper bound $(b-3q)/2$ on $p$; for
$2a-b\leq q\leq a/2$, use $a-2q$. The result is
\begin{align}\label{eq:volmiddle}
 V(a,b)=-\frac1{120}\bigl(&27a^6-108a^5b+180a^4b^2-160a^3b^3\notag\\
                         &+60a^2b^4-12ab^5+b^6\bigr).
\end{align}
These integrations involve only a polynomial of degree four, with linear
endpoints, so ordinary term-by-term integration suffices. The optional
symbolic companion verifies the three definite integrals exactly.

Finally, subtract \eqref{eq:volsmall} from \eqref{eq:volmiddle}; the difference
is $-(2b-3a)^6/3240$. Subtract \eqref{eq:volmiddle} from
\eqref{eq:vollarge}; the difference is $(b-2a)^6/120$. This gives
\eqref{eq:volumetruncated} in all three regions, including their shared
boundaries.
\end{proof}

For equal bounds, the first region applies and
\[
 V(a,a)=\frac{65}{648}a^6,
\]
agreeing with \eqref{eq:asymptotic}. The sixth-power corrections also show
that the leading volume and its derivatives through order five match across
each positive boundary ray; the change first becomes visible in sixth
transverse derivatives.

There is a direct volume interpretation. For continuous entries with bounds
$a,b$, fixed difference parameters leave four intervals of translation of
lengths $a-H_1,a-H_1,b-H_2,b-H_2$, when positive. Their product is the
four-dimensional fiber volume. Integrating over $x,y$, then using the twelve
chambers and the determinant-one change \eqref{eq:pq}, gives precisely
\eqref{eq:volumeintegral}. Hence $65/648$ is the volume of $\mathcal P$ in
the measure for which the integer lattice in its six parametrizing
coordinates has covolume one. It is \emph{not} the convention that multiplies
this volume by $6!$; under that alternative normalization the value is
$650/9$.

\section{Source audit: the shift and the printed constant}
\label{sec:source}

Let $T(n)$ denote the literal expression printed in Conjecture 11 of
Kauers and Koutschan \cite[p.~25]{kk}. Written schematically, it is
\begin{equation}\label{eq:printed}
 T(n)=\frac{130n^6+1560n^5+8125n^4+23400n^3
                 +u_rn^2+v_rn+w_r}{1296},
 \qquad r=n\bmod6,
\end{equation}
with the following printed lower coefficients:
\begin{center}
\begin{tabular}{crrr}
\toprule
$r$&$u_r$&$v_r$&$w_r$\\
\midrule
0&40788&42768&20736\\
1&40692&42128&20045\\
2&40788&42256&19712\\
3&40788&42768&20493\\
4&40692&42128&20288\\
5&40788&42256&19496\\
\bottomrule
\end{tabular}
\end{center}
Both the journal PDF and the arXiv version inspected for this report display
$19496$ in the final row.

There are two distinct observations. First, $T(0)=16$, while the entry-bound
count is $A(0)=1$. The beginning of the printed expression is shifted by one
relative to the entry bound. Second, the final row has a genuine arithmetic
inconsistency, not just a shift. At $n=5$,
\begin{equation}\label{eq:typo}
 T(5)=\frac{598513}{48}=12469+\frac1{48},
\end{equation}
which cannot be an array count.

Our exact formulas determine the comparison for every index:
\begin{equation}\label{eq:sourcecomparison}
 T(n)-A(n+1)=
 \begin{cases}
 1/48,&n\equiv5\pmod6,\\
 0,&\text{otherwise},
 \end{cases}
 \qquad n\geq0.
\end{equation}
Indeed, substitute $n=6k+r$ into \eqref{eq:printed} and compare coefficients
with Table~\ref{tab:residues}, using $A(6k+6)=A(6(k+1))$ for the final
class. Five polynomial differences vanish; the remaining numerator difference
is the constant $27$, giving $27/1296=1/48$.

Replacing $19496$ by $19469$ removes that constant difference. Thus the
corrected printed expression is exactly $A(n+1)$. Equivalently, substituting
$n-1$ into the corrected expression gives the count with entries bounded by
$n$, for $n\geq1$.

This audit is not the main solution. Finding the nonintegral value in
\eqref{eq:typo} would only refute a transcription-level statement. The
substantive result is the derivation from the array constraints of the
correct OEIS quasipolynomial, including every residue class and all indices.

\section{Verification, reproducibility, and computational scope}
\label{sec:verification}

The archive contains two verification programs. The main program,
\texttt{verify.py}, uses only the Python standard library and exact integers
or rational numbers. The optional \texttt{verify\_symbolic.py} uses SymPy for
independent differentiation, rational simplification, and polynomial
integration. Neither program requires network access.

\subsection{What was actually checked}
The main program completed thirteen groups of checks. Three are all-index
algebra certificates: the degree-eighteen identity
\eqref{eq:polycertificate}, the denominator-cleared partial fractions
\eqref{eq:PF}, and the degree-fifty-four identity
\eqref{eq:residuecertificate}. It also checks the noncancellation values
\eqref{eq:noncancel} and the numerator/denominator reversal identities.

For an independent combinatorial test, it constructs the three directions
from triangular-grid coordinates $(r,c)$, $0\leq c\leq r\leq4$, rather
than copying the parametrization into the validity test. It exhaustively
examines every filling for bounds zero, one, and two. The numbers of candidate
fillings are $1$, $4096$, and $531441$, and the numbers accepted are $1$,
$16$, and $105$. Each accepted filling is also sent through the inverse
parametrization and reconstructed.

The program additionally compares the full two-dimensional lattice sum with
the chamber sum for $0\leq n\leq100$, the chamber sum with the residue
formula for $0\leq n\leq300$, and the rational-function coefficients with
the residue formula for $0\leq n\leq1000$. All twenty-nine initial values
displayed on the retrieved OEIS page agree. It checks the auxiliary dihedral
orbits on $[-20,20]^2$, the centered formula and reciprocity for
$-200\leq n\leq200$, and three independent descriptions of $A(n,m)$ on
the rectangle $0\leq n,m\leq20$. The source comparison
\eqref{eq:sourcecomparison} is checked by exact polynomial composition,
not only by numerical substitution.

The optional symbolic program completed nine further exact checks: the two
Euler derivatives, the rational generating-function simplification, three
volume integrations, two boundary-correction identities, and the equal-bound
leading coefficient. The run used Python 3.13.5 and SymPy 1.14.0. Machine-readable
reports record the versions and measured run times.

\subsection{What these checks do and do not prove}
The manuscript proves the integer bijection, translation count, fundamental
domain, and summation identity for all indices. Those steps are the
mathematical foundation. The polynomial certificates verify finite algebra
identities that themselves imply all-index generating-function statements.
The finite enumerations serve a different purpose: they independently detect
mistakes in the interpretation of the array constraints, indexing, chamber
weights, or copied coefficients.

A finite test is not promoted to a general proof merely because its range is
large. Conversely, a finite \emph{coefficient comparison between two explicitly
bounded-degree polynomials} proves their identity. The verifier keeps these
two types of evidence separate. None of the files constitutes a
proof-assistant formalization, and the numerical tests cannot by themselves
establish literature priority.

\subsection{Rebuilding and using the artifacts}
From the extracted archive directory, run
\begin{verbatim}
python3 verify.py
python3 verify_symbolic.py      # optional; requires SymPy
pdflatex -interaction=nonstopmode -halt-on-error article.tex
pdflatex -interaction=nonstopmode -halt-on-error article.tex
\end{verbatim}
The second LaTeX pass resolves the contents and cross-references.
\texttt{data/sequence.csv} contains $A(n)$ for $0\leq n\leq1000$;
\texttt{data/two\_bound\_counts.csv} contains the checked two-bound rectangle.
\texttt{README.md} documents the files, conventions, and status. The supplied
\texttt{Makefile} provides equivalent build targets.

For evaluation at a single very large nonnegative index, the residue
polynomial is the most economical formula: choose the residue and use Horner's
rule in $k=\lfloor n/6\rfloor$. It uses a fixed number of integer arithmetic
operations, although their bit cost grows with the size of $n$. The chamber
sum is preferable as a transparent independent enumerator, with $O(n^2)$
terms. The order-fourteen recurrence is convenient for generating a long
prefix. These are arithmetic-operation statements, not claims of constant
bit complexity.

\section{Conclusion and remaining directions}

The explicit A195806 quasipolynomial can be derived from an elementary
structure that is not apparent in the twelve-entry definition. Four
translation variables account for the fourth-degree weight, while two
difference variables account for the sixth-degree total growth. The two
translation ranges are hexagonal norms with a common twelvefold parameter
symmetry. In a fundamental chamber their values become $p+2q$ and $2p+3q$,
which explains the appearance of the denominators $1-z^2$ and $1-z^3$ and,
after cancellation, the exact quasiperiod six.

The main result is a proof of the correctly indexed formulas and recurrence
still marked conjectural in the source inspected, together with minimality,
reciprocity, and a two-bound extension. It does not resolve every sequence in
the 2023 paper, establish a new general lattice-point algorithm, or prove that
no earlier solution to this particular conjecture exists. Those are separate
claims and are not needed here.

A natural next investigation is the analogous triangular-array problem at
other side lengths. The successful part of the present method is not simply
linear elimination, which remains available in general, but the subsequent
discovery of repeated translation ranges and a common manageable chamber
structure. Determining when those features persist is a concrete structural
question. No unverified general formula for larger triangles is asserted in
this report.

\appendix
\section{Coefficient lists and a compact counting implementation}
\label{app:implementation}

For use in independent implementations, the denominator expands as
\begin{align*}
 Q(z)={}&1-3z+2z^2-z^3+6z^4-5z^5-3z^6\notag\\
       &\quad+3z^8+5z^9-6z^{10}+z^{11}-2z^{12}+3z^{13}-z^{14}.
\end{align*}
The numerator and denominator coefficients, in ascending order, are
\begin{verbatim}
P = [1,13,59,212,471,753,882,753,471,212,59,13,1]
Q = [1,-3,2,-1,6,-5,-3,0,3,5,-6,1,-2,3,-1]
\end{verbatim}
Since $Q(0)=1$, coefficients of $P/Q$ can be generated recursively as
\[
 A(n)=P_n-\sum_{j=1}^{\min(n,14)}Q_jA(n-j),
\]
where $P_n=0$ for $n>12$. This formula incorporates the initial terms and
avoids a separate recurrence initialization convention.

The core chamber enumerator is short enough to reproduce in full:
\begin{verbatim}
def count(n: int) -> int:
    if isinstance(n, bool) or not isinstance(n, int) or n < 0:
        raise ValueError("n must be a nonnegative integer")
    total = 0
    for p in range(n // 2 + 1):
        for q in range((n - 2*p) // 3 + 1):
            r = n - 2*p - 3*q
            mu = 1 if p == q == 0 else 6 if p*q == 0 else 12
            total += mu * (r+1)**2 * (r+p+q+1)**2
    return total
\end{verbatim}
Every term is nonnegative, and no floating-point operation is involved.
The complete verifier adds error checks, independent enumerators, exact
polynomial arithmetic, reports, and data export.

\section{A short numerical reference table}
\label{app:values}

\begin{center}
\begin{tabular}{rr@{\qquad}rr@{\qquad}rr}
\toprule
$n$&$A(n)$&$n$&$A(n)$&$n$&$A(n)$\\
\midrule
0&1&7&27412&14&1155567\\
1&16&8&55059&15&1699684\\
2&105&9&102952&16&2442553\\
3&496&10&181543&17&3438468\\
4&1759&11&304908&18&4752283\\
5&5052&12&491563&19&6460432\\
6&12469&13&765184&20&8652429\\
\bottomrule
\end{tabular}
\end{center}
These values are included for convenient indexing checks. The exact proof
is provided by the preceding bijection, chamber reduction, and generating
function, not by this table.

\begin{thebibliography}{9}

\bibitem{oeis}
R. H. Hardin, \emph{A195806: Number of triangular size-five arrays with
entries $0,\ldots,n$, equal sums on lines of equal length, and zero corners},
The On-Line Encyclopedia of Integer Sequences, introduced September 23, 2011.
Conjectured recurrence and quasipolynomial contributed by M. Kauers and
C. Koutschan, March 1, 2023.
\url{https://oeis.org/A195806}.
Retrieved September 20, 2026.

\bibitem{kk}
M. Kauers and C. Koutschan, \emph{Some D-Finite and Some Possibly D-Finite
Sequences in the OEIS}, Journal of Integer Sequences \textbf{26} (2023),
Article 23.4.5, 42 pp. See Section 5.1, especially Conjecture 11 on p.~25.
\url{https://cs.uwaterloo.ca/journals/JIS/VOL26/Koutschan/kout4.pdf}.
Also arXiv:2303.02793, version 2, April 24, 2023,
\url{https://arxiv.org/abs/2303.02793}.
Journal and arXiv versions inspected September 20, 2026.

\end{thebibliography}
\end{document}
